Before we talk about SNARKs (specifically, PLONK), it helps to separate
out an ingredient that underlies much of programmable cryptography,
which is the idea of a \emph{polynomial commitment}. Specifically, we
will talk about the KZG polynomial commitment, which plays an important
role in PLONK (and many other protocols). For a higher-resolution
understanding of KZG, it helps to understand \emph{elliptic curves}
(especially in the context of \emph{pairings}), which are ubiquitous in
cryptography. If you are uninterested (or experienced) in mathematical
details, you can and should skip elliptic curves and jump to Section \ref{kzg}.
If you are comfortable with black-boxing Section \ref{kzg}, you can even jump
straight to SNARKs into the next chapter.

The roadmap goes roughly as follows:

\begin{itemize}
\item In Section \ref{ec} we will define \alert{elliptic curves} and
  describe one standard elliptic curve $E$, the \emph{BN254 curve},
  that will be used in these notes.

\item In Section \ref{discretelog} we describe the
  \alert{discrete logarithm assumption} (Assumption \ref{tcb:ddh}),
  which we need to make to provide security to our protocols.
  As an example, in Section \ref{eddsa} we describe how
  Assumption \ref{tcb:ddh} can be used to construct a signature scheme,
  namely EdDSA.\footurl{https://en.wikipedia.org/wiki/EdDSA}

\item The EdDSA idea will later grow up to be the KZG commitment scheme in
  Section \ref{kzg}.
\end{itemize}

\section{Elliptic curves}\label{ec}
Every modern cryptosystem rests on a hard problem -- a computationally
infeasible challenge whose difficulty makes the protocol secure. The
best-known example is RSA,\footurl{https://en.wikipedia.org/wiki/RSA\_(cryptosystem)}
which is secure because it is hard to factor a composite number (like
$6887$) into prime factors ($6887 = 71 \cdot 97$).

Our SNARK protocol will be based on a different hard problem: the
``discrete logarithm problem''
(see Section \ref{discretelog}) on elliptic curves. But
before we get to the problem, we need to introduce some of the math
behind elliptic curves.

An \alert{elliptic curve} is a set of points with a group operation. The
set of points is the set of solutions $(x,y)$ to an equation in two
variables; the group operation is a rule for ``adding'' two of the
points to get a third point. Our first task will be to understand what
all this means. Rather than set up a general definition of elliptic
curve, for these notes we will be satisfied to describe one specific
elliptic curve that can be used for all the protocols we describe later.
The curve we choose for these notes is the \emph{BN254 curve}.

The BN254 specification fixes a specific\footnote{If you must know, the
  values in the specification are given exactly by
  \begin{align*}
  x & \defeq 4965661367192848881 \\
  p & \defeq 36x^{4} + 36x^{3} + 24x^{2} + 6x + 1 \\
  & = 218882428718392752222464057452572750886 96311157297823662689037894645226208583 \\
  q & \defeq 36x^{4} + 36x^{3} + 18x^{2} + 6x + 1 \\
  & = 218882428718392752222464057452572750885 48364400416034343698204186575808495617.
  \end{align*}
}
large prime $p \approx 2^{254}$ (and a second large
prime $q \approx 2^{254}$ that we define later) which has been
specifically engineered to have certain properties (Jonathan Wang has a
blog post\footurl{https://hackmd.io/@jpw/bn254} about the properties of
this curve). The name BN stands for Barreto-Naehrig, two mathematicians
who proposed a family of such curves in 2006.\footurl{https://link.springer.com/content/pdf/10.1007/11693383\_22.pdf}

\begin{definition}{BN254 curve}{}
The \emph{BN254 curve} is the elliptic curve
\begin{equation}
\label{bn254eqn}
E(\FF_p) : Y^2 = X^3 + 3
\end{equation}
defined over $\FF_p$, where $p \approx 2^{254}$
is the prime from the BN254 specification.
\end{definition}

So each point on $E( \FF_{p} )$ is an ordered pair
$(X,Y) \in \FF_{p}^{2}$ satisfying Equation \eqref{bn254eqn}.
Okay, actually, that is a white lie: Conventionally, there is one additional
point $O = (0,\infty)$ called the ``point at infinity'' added in
(whose purpose we describe in the next section).

The constants $p$ and $q$ are contrived so that the following holds:

\begin{theorem}{BN254 has prime order}{}
  Let $E$ be the BN254 curve. The number of
  points in $E( \FF_{p} )$, including the point at
  infinity $O$, is a prime $q \approx 2^{254}$.
\end{theorem}

\begin{definition}{Baby Jubjub prime}{}
This prime $q \approx 2^{254}$ is affectionately called the \alert{Baby Jubjub prime}.
It will usually be denoted by $q$ in these notes.
\end{definition}
(The name ``Baby Jubjub'' is a reference to
\emph{The Hunting of the Snark}\footurl{https://en.wikipedia.org/wiki/The\_Hunting\_of\_the\_Snark}.)

So, at this point, we have a bag of $q$ points denoted
$E( \FF_{p} )$.
However, right now it only has the structure of a set.

The beauty of elliptic curves is that it is possible to define an
\textbf{addition} operation on the curve; this is called the \alert{group law}
on the elliptic curve.\footurl{https://en.wikipedia.org/wiki/Elliptic\_curve\#The\_group\_law}
This addition will make $E( \FF_{p} )$ into an
abelian group whose identity element is the point at infinity $O$.

This group law involves some heavy algebra. It is not important to
understand exactly how it works. All you really need to take away from
this section is that there is some group law, and we can program a
computer to compute it. We provide details in Appendix~\ref{appendix-ec-mult}
for the interested reader.

In summary, we have endowed the set of points
$E( \FF_{p} )$ with the additional structure of an
abelian group, which happens to have exactly $q$ elements. However, an
abelian group with prime order is necessarily cyclic. In other words:

\begin{theorem}{The group BN254 is isomorphic to $\boldmath{\ZZ/q\ZZ}$}{}
  Let $E$ be the BN254 curve.
  We have the isomorphism of abelian groups
  \[E( \FF_{p} ) \cong \ZZ/q\ZZ.\]
\end{theorem}

In these notes, this isomorphism will basically be a standing assumption.
Moving forward, we will abuse notation slightly and just
write $E$ instead of $E( \FF_{p} )$.
In fancy language, $E$ will be a one-dimensional vector space over $\FF_{q}$.
In less fancy language, we will be working with points on $E$ as black boxes.
We will be able to add them, subtract them, and multiply them by arbitrary scalars from $\FF_{q}$.

Consequently --- and this is important --- \textbf{one should think of
$\FF_{q}$ as the base field for all our cryptographic primitives}
(despite the fact that the coordinates of our points are in $\FF_{p}$).

\begin{remark}{}{}
Whenever we talk about protocols, and there are any sorts of
``numbers'' or ``scalar'' in the protocol,
these scalars are always going to be elements of $\FF_q$.
Since $q \approx 2^{254}$,
that means we are doing something like $256$-bit integer arithmetic.
This is why the baby Jubjub prime $q$ gets a special name,
while the prime $p$ is unnamed and does not get any screen-time later.
\end{remark}

\section{Discrete logarithm}\label{discretelog}

For our systems to be useful, rather than relying on factoring, we will
rely on the so-called \alert{discrete logarithm assumption}.

\begin{assumption}{Discrete logarithm assumption}{ddh}
Let $E$ be the BN254 curve (or another
standardized curve). Given arbitrary nonzero $g,g' \in E$, the
\alert{discrete logarithm problem} asks you to find an integer $n$ such
that $n \cdot g = g'$.

Experience suggests that the discrete logarithm problem is hard:
In general, we do not know a fast algorithm to solve it.
The \alert{discrete logarithm assumption} says that no such algorithm exists.
\end{assumption}

In other words, if one only sees $g \in E$ and $n \cdot g \in E$,
one cannot find $n$. For cryptography, we generally assume $g$ has
order $q$, so we will talk about $n \in {\mathbb{N}}$ and
$n \in \FF_{q}$ interchangeably. In other words, $n$ will
generally be thought of as being up to about $2^{254}$ in size.

\begin{remark}{The name ``discrete log''}{}
This problem is called ``discrete log'' because if one uses multiplicative notation
for the group operation, it looks like solving $g^n = g'$ instead.
We will never use this multiplicative notation in these notes.
\end{remark}

On the other hand, given $g \in E$, one can compute $n \cdot g$ in
just $O( \log n )$ operations, by repeated
squaring.\footurl{https://en.wikipedia.org/wiki/Exponentiation\_by\_squaring}
For example, to compute $400g$, one only needs to do $10$ additions,
rather than $400$: One starts with \[\begin{aligned}
2g & = g + g \\
4g & = 2g + 2g \\
8g & = 4g + 4g \\
16g & = 8g + 8g \\
32g & = 16g + 16g \\
64g & = 32g + 32g \\
128g & = 64g + 64g \\
256g & = 128g + 128g
\end{aligned}\] and then computes \[400g = 256g + 128g + 16g.\]

Because we think of $n$ as up to $q \approx 2^{254}$-ish in size, we
consider $O( \log n )$ operations like this to be quite
tolerable.

\subsection{Curves other than BN254}
We comment briefly on how the previous definitions adapt to other
curves, although readers could get away with always assuming $E$ is
BN254 if they prefer.

In general, we could have chosen for $E$ any equation of the form
$Y^{2} = X^{3} + aX + b$ and chosen any prime $p \geq 5$ such that a
nondegeneracy constraint $4a^{3} + 27b^{2} \not \equiv 0 \pmod p$ holds.
In such a situation, $E( \FF_{p} )$ will
indeed be an abelian group once the identity element $O = (0,\infty)$ is added in.

How large is $E( \FF_{p} )$?
There is a theorem called Hasse's theorem\footurl{https://en.wikipedia.org/wiki/Hasse's_theorem_on_elliptic_curves}
that states the number of points in $E( \FF_{p} )$
is between $p + 1 - 2\sqrt{p}$ and $p + 1 + 2\sqrt{p}$. But there is
no promise that $E( \FF_{p} )$ will be \emph{prime};
consequently, it may not be a cyclic group either.
So among many other considerations, the choice of constants in BN254 is
engineered to get a prime order.

There are other curves used in practice for which
$E( \FF_{p} )$ is not a prime,
but rather a small multiple of a prime.
The popular Curve25519\footurl{https://en.wikipedia.org/wiki/Curve25519}
is such a curve that is also believed to satisfy Assumption \ref{tcb:ddh}.
Curve25519 is defined as
\[ Y^{2} = X^{3} + 486662X^{2} + X \]
over $\FF_{p}$ for the
prime $p \defeq 2^{255} - 19$. Its order is $8$ times a large prime
\[q' \defeq 2^{252} + 27742317777372353535851937790883648493. \]
In that case, to generate a random point on Curve25519 with order $q'$,
one will usually take a random point on the curve and multiply it by $8$.

BN254 is also engineered to have a property called
\emph{pairing-friendliness}, which is defined in Section \ref{pairing-friendly} when
we need it later. (In contrast, Curve25519 does not have this property.)

\subsection{Example application: EdDSA signature scheme}\label{eddsa}
We will show how Assumption \ref{tcb:ddh} can be used to construct a
signature scheme that replaces RSA.
This scheme is called EdDSA,\footurl{https://en.wikipedia.org/wiki/EdDSA}
and it is used quite frequently (e.g., in OpenSSH and GnuPG).
One advantage it has over RSA is that its key size is much smaller:
both the public and private key are 256 bits.
(In contrast, RSA needs 2048-4096 bit keys for comparable security.)

\begin{definition}{}{armor}
Let $E$ be an elliptic curve and let $g \in E$
be a fixed point on it of prime order $q \approx 2^{254}$.
For $n \in \ZZ$ (equivalently $n \in \FF_q$) we define
\[ [n] \defeq n \cdot g \in E. \]
\end{definition}

The hardness of the discrete logarithm means that, given
$[n]$, we cannot get $n$. You can almost think of the
notation as an ``armor'' on the integer $n$: It conceals the integer,
but still allows us to perform (armored) addition:
\[ [a + b] = [a] + [b]. \] In other
words, $n \mapsto [n]$ viewed as a map
$\FF_{q} \rightarrow E$ is $\FF_{q}$-linear.

So now suppose Alice wants to set up a signature scheme.

\begin{algo}{EdDSA public and secret key}{}
\begin{enumerate}
\item
  Alice picks a random integer $d \in \FF_{q}$ as her
  \emph{secret key} (a piece of information that she needs to keep
  private for the security of the protocol).
\item
  Alice publishes $[d] \in E$ as her \emph{public key} (a
  piece of information which, even when obtained by adversaries, does
  not challenge the security of the protocol).
\end{enumerate}
\end{algo}

Now suppose Alice wants to prove to Bob that she approves the message
$\textsf{msg}$, given her published public key
$[d]$.

\begin{algo}{EdDSA signature generation}{}
Suppose Alice wants to sign a message $\textsf{msg}$.

\begin{enumerate}
\item
  Alice picks a random scalar $r \in \FF_{q}$ (keeping this
  secret) and publishes $[r] \in E$.
\item
  Alice generates a number $n \in \FF_{q}$ by hashing
  $\textsf{msg}$ with all public information, say
  \[ n \defeq \hash([r],\textsf{msg},[d]). \]
\item
  Alice publishes the integer \[ s \defeq (r + dn) \bmod q. \]
\end{enumerate}

In other words, the signature is the ordered pair
$( [r],s )$.
\end{algo}

\begin{algo}{EdDSA signature verification}{}
For Bob to verify a signature
$( [r],s )$ for $\textsf{msg}$:

\begin{enumerate}
\item
  Bob recomputes $n$ (by also performing the hash) and computes
  $[s] \in E$.
\item
  Bob verifies that
  $[r] + n \cdot [d] = [s]$.
\end{enumerate}
\end{algo}

An adversary cannot forge the signature even if they know $r$ and
$n$. Indeed, such an adversary can compute what the point
$[s] = [r] + n[d]$ should be,
but without knowledge of $d$ they cannot get the integer $s$, due to
Assumption \ref{tcb:ddh}.

The number $r$ is called a \emph{blinding factor} because its use
prevents Bob from stealing Alice's secret key $d$ from the published
$s$. It is therefore imperative that $r$ is not known to Bob nor
reused between signatures, and so on. One way to do this would be to
pick $r = \hash(d,\textsf{msg})$; this has the
bonus that it is deterministic as a function of the message and signer.

In Section \ref{kzg} we will use ideas quite similar to this to build the KZG
commitment scheme.

\subsection{Example application: Pedersen commitments}
\label{pedersen}

A \alert{commitment scheme} is a protocol where Alice wants to commit
some value $x$ to Bob that is later revealed. Typically Alice gives
Bob some ``commitment'' $\Com(x)$ and later reveals
$x$. What we want is that this protocol is both \emph{binding} (Alice
cannot change her mind about $x$ depending on Bob's later actions) and
\emph{hiding} (Bob does not get any information about $x$ from
$\Com(x)$). The KZG scheme we are building towards will
be a commitment scheme for polynomials, but we can already use elliptic
curves to commit \textbf{numbers} with something called a Pedersen
commitment, which we will now describe.

A multivariable generalization of Assumption \ref{tcb:ddh} is that if
$g_{1},\ldots,g_{n} \in E$ are a bunch of randomly chosen points of
$E$ with order $q$, then it is computationally infeasible to find
$( a_{1},\ldots,a_{n} ) \neq ( b_{1},\ldots,b_{n} ) \in \FF_{q}^{n}$
such that
\[a_{1}g_{1} + \cdots + a_{n}g_{n} = b_{1}g_{1} + \cdots + b_{n}g_{n}.\]
(Remember that $q \approx 2^{256}$ is very large.)

\begin{definition}{Computational basis}{comp}
In these notes, if there is a globally known elliptic curve $E$
and points $g_1, \ldots, g_n$ have order $q$ and no known nontrivial
linear dependencies between them,
we will say they are a \alert{computational basis over $\FF_q$}.
\end{definition}

\begin{remark}{}{}
This may horrify pure mathematicians because we are pretending the map
\[ \FF_q^n \rightarrow \FF_q \text{ by } (a_1, \ldots, a_n) \mapsto \sum_1^n a_i g_i \]
is injective,
even though the domain is an $n$-dimensional $\FF_q$-vector space
and the codomain is one-dimensional.
This can feel weird because our instincts from linear algebra in pure math
are wrong now. This map, while not injective in theory,
ends up being injective \emph{in practice} (because we cannot find collisions).
And this is a critical standing assumption for this entire framework!
\end{remark}

This injectivity gives us a sort of hash function on vectors (with
``linearly independent'' now being phrased as ``we cannot find a
collision''). To spell this out:

\begin{definition}{Pedersen comittment}{}
Let $g_1, \ldots, g_n \in E$ be a computational basis over $\FF_q$.
Given a vector
\[ \vec{a} = \langle a_1, \ldots, a_n \rangle \in \FF_q^n \] of scalars,
the group element
\[ \sum a_i g_i = a_1 g_1 + \cdots + a_n g_n \in E \]
is called the \alert{Pedersen commitment} to our vector $\vec{a}$.
\end{definition}

The Pedersen commitment is thus a sort of hash function: Given the group
element above, one cannot recover any of the $a_{i}$; but given the
entire vector $\overrightarrow{a}$ one can compute the Pedersen
commitment easily.

We will not use Pedersen commitments in this book, but they will be
closely related to KZG.

\subsection{Appendix: The group law of an elliptic curve}
\label{appendix-ec-mult}

As promised in Section~\ref{ec},
this optional appendix describes the details of the math
behind the group law of the elliptic curve.

So, let us get started. The equation of $E$ is cubic -- the
highest-degree terms have degree $3$. This means that (in general) if
you take a line $y = mx + b$ and intersect it with $E$, the line
will meet $E$ in exactly three points. The basic idea behind the group
law is: If $P,Q,R$ are the three intersection points of a line (any
line) with the curve $E$, then the group-law addition of the three
points is \[P + Q + R = O.\]

(You might be wondering how we can do geometry when the coordinates
$x$ and $y$ are in a finite field. It turns out that all the
geometric operations we are describing -- like finding the intersection
of a curve with a line -- can be translated into algebra, and then you
just do the algebra in your finite field. We will come back to this.)

Why three points? Algebraically, if you take the equations
$Y^{2} = X^{3} + 3$ and $Y = mX + b$ and try to solve them, you get
\[(mX + b)^{2} = X^{3} + 3,\] which is a degree-3 polynomial in $X$,
so it has (at most) three roots. In fact if it has two roots, it is
guaranteed to have a third (because you can factor out the first two
roots, and then you are left with a linear factor).

Now given two points $P$ and $Q$, how do we find their sum
$P + Q$? We can draw the line through the two points. That line --
like any line -- will intersect $E$ in three points: $P$, $Q$, and
a third point $R$. Now since $P + Q + R = 0$, we know that
\[- R = P + Q.\]

So now the question is just: how to find $- R$? Well, it turns out
that if $R = ( x_{R},y_{R} )$, then
\[- R = ( x_{R}, - y_{R} ).\] Why is this? If you take the
vertical line $X = x_{R}$, and try to intersect it with the curve, it
looks like there are only two intersection points. After all, we are
solving \[Y^{2} = x_{R}^{3} + 3,\] and since $x_{R}$ is fixed now,
this equation is quadratic. The two roots are $Y = \pm y_{R}$.

There are only two intersection points, but we say that the third
intersection point is ``the point at infinity'' $O$. (The reason for
this lies in projective geometry, but we will not get into it.) So the
group law here tells us
\[( x_{R},y_{R} ) + ( x_{R}, - y_{R} ) + O = O.\]
And since $O$ is the identity, we get
\[- R = ( x_{R}, - y_{R} ).\]

So:

\begin{itemize}
\item Given a point $P = ( x_{P},y_{P} )$, its negative is just
  $- P = ( x_{P}, - y_{P} )$.
\item To add two points $P$ and $Q$, compute the line through the two
  points, let $R$ be the third intersection of that line with $E$,
  and set \[P + Q = - R.\]
\end{itemize}

We just described the group law as a geometric thing, but there are
algebraic formulas to compute it as well. They are kind of a mess, but
here goes.

If $P = ( x_{P},y_{P} )$ and
$Q = ( x_{Q},y_{Q} )$, then the line between the two points
is $Y = mX + b$, where \[m = \frac{y_{Q} - y_{P}}{x_{Q} - x_{P}}\] and
\[b = y_{P} - mx_{P}.\]

The third intersection is $R = ( x_{R},y_{R} )$, where
\[x_{R} = m^{2} - x_{P} - x_{Q}\] and \[y_{R} = mx_{R} + b.\]

There are separate formulas to deal with various special cases (if
$P = Q$, you need to compute the tangent line to $E$ at $P$, for
example), but we will not get into it.
